\section{The interpolation upper bound}\label{upper:section}

The interpolation method of Guerra~\cite{Guerra03}, adapted to diluted
systems by Franz and Leone~\cite{FL03} and to arbitrary finite levels by
Panchenko and Talagrand~\cite{PT04}, replaces each real interaction by
$p$ independent cavity interactions. The difference is
controlled by the convexity of $x\mapsto x^p$. We give the argument for
the full trial class, including unbounded messages and integrable
physical disorder.

\begin{proposition}\label[proposition]{upper:bound}
Under the hypotheses of \cref{main:theorem},
$F_N\le\cB_r(\zeta,\mathbf m)$ for every $N$ and every admissible finite-level
trial law. Consequently $F_N\le\inf_{r\ge0}\Phi_r$.
\end{proposition}

\begin{proof}
First dispose of the degenerate factorization. Put
$M=\max_{s=\pm1}|f(s)|$. If $b=0$ almost surely or $M=0$ almost surely,
the interaction is a spin-independent random constant $c$, and
\[
 F_N=\cB_r=\log2+\E\log\cosh h+\alpha\E c
\]
for every trial law. Otherwise choose $u,v>0$ with
$\P(M\ge u)>0$ and $\P(|b|\ge v)>0$. Since
$|b|\prod_{\ell=1}^p M_\ell<1$ almost surely, independence gives
\[
 \|b\|_\infty\le u^{-p},\qquad
 \|M\|_\infty\le(vu^{p-1})^{-1}.
\]
Indeed, violating either bound on a set of positive probability and
restricting the other independent factors to these lower bounds would
violate the strict product condition. All moments used below are
therefore finite.

\emph{Interpolation.}
Fix a trial law and set $L=r+1$. Interpolate between real interactions of mean $t\alpha N$
and, at each site $i$, independent cavity interactions
$U_\theta(x_1,\ldots,x_{p-1};\sigma_i)$ of mean $(1-t)\alpha p$.
Different cavity terms have disjoint message labels. Counts, interactions,
fields, and root message measures are quenched. Apply the backward
power-mean recursion to the spin partition function and denote its expected
root logarithm divided by $N$ by $\varphi(t)$. For $r=0$ all messages are
quenched and there is no auxiliary recursion. Independence between different
sites throughout the recursion gives
\[
 \varphi(1)=F_N,\qquad
 \varphi(0)=\log2+\E\log T_0V_{\rm s}.
\]

Poisson differentiation gives $\varphi'(t)/\alpha$ as the expected insertion
increment of one real interaction minus $p$ times the increment of one
cavity interaction at a uniform site. Each increment has absolute value at
most $B(\theta)$, since the recursion preserves multiplicative bounds.
Conditioning on the counts justifies the derivative on $(0,1)$, and the same
bound gives continuity at the endpoints.

We claim that
\begin{equation}\label{upper:derivative}
 \varphi'(t)+\alpha(p-1)\E\log T_0V_{\rm e}\le0.
\end{equation}
All three insertions have the form $a(1+D)$ under
\eqref{model:factorization}. For a real interaction,
$D=b\prod_\ell f_\ell(\sigma_{i_\ell})$; for a cavity interaction, replace
$p-1$ of the factors by
$\bar f(x)=\sum_s f(s)q_x(s)$; for an edge insertion, replace all $p$.
In the last case the independent message process is appended to the
unperturbed tilted tree, so its increment is $\log T_0V_{\rm e}$.

\emph{Normalization.}
We remove $a$ before taking expectations, so no integrability of
$\log a$ is needed. Put
$D_*=b\prod_\ell f_\ell(1)$ and $c=\theta(1,\ldots,1)$.
Subtract $c$ from each insertion increment. These integrable subtractions
cancel in the combination with coefficients $1,-p,p-1$.
Keep the interpolated system fixed. For $0<\lambda<1$, replace
each normalized increment by
\begin{equation}\label{upper:regularization}
 \log Y_0(1+\lambda D)-\log(1+\lambda D_*),
\end{equation}
where $Y_0$ is the insertion recursion of \cref{tree:log-expansion}.
For every spin input, the ratio
\[
 \frac{1+\lambda D}{1+\lambda D_*}
 =(1-w)+w\frac{1+D}{1+D_*},
 \qquad w=\frac{\lambda(1+D_*)}{1+\lambda D_*}\in[0,1],
\]
lies between $e^{-2B(\theta)}$ and $e^{2B(\theta)}$. The same is true
after averaging any subset of the spin inputs against the $q_x$'s.
It follows that \eqref{upper:regularization} is bounded in absolute value by
$2B(\theta)$. For each fixed interaction, convergence as $\lambda\uparrow1$
is uniform in the spin inputs and messages and is preserved by the finite
recursion. Dominated convergence therefore recovers the desired normalized
increments.

\emph{Convexity.}
Apply \cref{tree:log-expansion} to the regularized increments. On a tree
$\cT$ with $n$ replica leaves define
\[
 P_{\cT}=\frac1N\sum_{i=1}^N\E_f\prod_{a=1}^n f(\sigma_i^a),\qquad
 B_{\cT}=\E_{f,x}\prod_{a=1}^n\bar f(x^a).
\]
Here $\E_f$ averages one fresh function $f$, shared by all $n$ replica
factors in the product. The expectation $\E_{f,x}$ also averages an
entire fresh message-label tree, including its root. This function and
message tree are independent of one another and of the old tilted system.
Thus $B_{\cT}$ is deterministic for a fixed tree, while $P_{\cT}$ still
depends on the old spin replicas. At $r=0$, the message itself is a root
variable shared by all replica leaves. Averaging the
independent uniform indices and the independent copies of $f$ and of
the new message label gives, respectively,
$P_{\cT}^p$, $P_{\cT}B_{\cT}^{p-1}$, and $B_{\cT}^p$.
The inserted $b$ is independent of these variables. Thus the order-$n$
coefficient of the combined regularized increment is
\begin{equation}\label{upper:remainder}
 -\frac{\lambda^n}{n}\E[(-b)^n]\,
 \E_{Q_n}\E\left[
 P_{\cT}^p-pP_{\cT}B_{\cT}^{p-1}+(p-1)B_{\cT}^p\right].
\end{equation}
The regularized baseline terms cancel. The tree law $Q_n$ is that of
\cref{tree:log-expansion}; it is common to the three insertions.
The series is absolutely convergent for $\lambda<1$, with each insertion
controlled by $\sum_{n\ge1}\lambda^n/n$. Since $p$ is even,
convexity of $x\mapsto x^p$ makes the bracket nonnegative, and
\eqref{model:positivity} makes \eqref{upper:remainder} nonpositive.
Taking $\lambda\uparrow1$ proves \eqref{upper:derivative}. Integration over
$t\in[0,1]$ gives the proposition.
\end{proof}
